\section{Complete opposite width-four encoder base trace}
\label{app:width4-opposite-trace}
This appendix discharges the finite base case used by the opposite shell induction.  Vertices and edge indices use the convention of \cref{app:encoders}.  A face $(a;x,y,z)$ is the unit square with fixed coordinate axis $a$ and lower corner $(x,y,z)$; $O$ denotes the nonsquare outer face.  Each row below is an explicit incidence check, so no program output is needed to infer connectivity.
The opposite base contains 80 vertices and the 79 selected tree edges printed in \cref{app:encoders}.  Starting at $(0,0,0)$, the following table gives one parent edge for every other vertex.
\begin{longtable}{@{}rrr@{}}
\toprule child & parent & edge index\\ \midrule
\endfirsthead \toprule child & parent & edge index\\ \midrule \endhead
$(0,1,0)$ & $(0,0,0)$ & $1$\\
$(1,0,0)$ & $(0,0,0)$ & $2$\\
$(0,1,1)$ & $(0,1,0)$ & $11$\\
$(2,0,0)$ & $(1,0,0)$ & $42$\\
$(0,0,1)$ & $(0,1,1)$ & $4$\\
$(0,2,1)$ & $(0,1,1)$ & $15$\\
$(1,1,1)$ & $(0,1,1)$ & $16$\\
$(3,0,0)$ & $(2,0,0)$ & $82$\\
$(1,0,1)$ & $(0,0,1)$ & $5$\\
$(0,3,1)$ & $(0,2,1)$ & $26$\\
$(1,2,1)$ & $(0,2,1)$ & $27$\\
$(2,1,1)$ & $(1,1,1)$ & $56$\\
$(4,0,0)$ & $(3,0,0)$ & $122$\\
$(1,0,2)$ & $(1,0,1)$ & $43$\\
$(2,0,1)$ & $(1,0,1)$ & $45$\\
$(1,3,1)$ & $(0,3,1)$ & $36$\\
$(1,2,2)$ & $(1,2,1)$ & $65$\\
$(2,2,1)$ & $(1,2,1)$ & $67$\\
$(2,1,0)$ & $(2,1,1)$ & $91$\\
$(3,1,1)$ & $(2,1,1)$ & $96$\\
$(0,0,2)$ & $(1,0,2)$ & $8$\\
$(2,0,2)$ & $(1,0,2)$ & $48$\\
$(3,0,1)$ & $(2,0,1)$ & $85$\\
$(2,3,1)$ & $(1,3,1)$ & $76$\\
$(0,2,2)$ & $(1,2,2)$ & $30$\\
$(2,2,2)$ & $(1,2,2)$ & $70$\\
$(3,2,1)$ & $(2,2,1)$ & $107$\\
$(1,1,0)$ & $(2,1,0)$ & $53$\\
$(3,1,0)$ & $(2,1,0)$ & $93$\\
$(0,1,2)$ & $(0,0,2)$ & $7$\\
$(3,0,2)$ & $(2,0,2)$ & $88$\\
$(4,0,1)$ & $(3,0,1)$ & $125$\\
$(3,3,1)$ & $(2,3,1)$ & $116$\\
$(0,3,2)$ & $(0,2,2)$ & $29$\\
$(2,3,2)$ & $(2,2,2)$ & $109$\\
$(3,2,2)$ & $(2,2,2)$ & $110$\\
$(1,2,0)$ & $(1,1,0)$ & $52$\\
$(4,1,0)$ & $(3,1,0)$ & $133$\\
$(0,1,3)$ & $(0,1,2)$ & $17$\\
$(1,1,2)$ & $(0,1,2)$ & $19$\\
$(4,0,2)$ & $(4,0,1)$ & $162$\\
$(4,1,1)$ & $(4,0,1)$ & $163$\\
$(0,3,3)$ & $(0,3,2)$ & $37$\\
$(1,3,2)$ & $(2,3,2)$ & $78$\\
$(3,3,2)$ & $(2,3,2)$ & $118$\\
$(0,2,0)$ & $(1,2,0)$ & $24$\\
$(2,2,0)$ & $(1,2,0)$ & $64$\\
$(4,2,0)$ & $(4,1,0)$ & $168$\\
$(0,0,3)$ & $(0,1,3)$ & $9$\\
$(2,1,2)$ & $(1,1,2)$ & $59$\\
$(1,3,3)$ & $(0,3,3)$ & $39$\\
$(4,3,2)$ & $(3,3,2)$ & $158$\\
$(0,3,0)$ & $(0,2,0)$ & $23$\\
$(3,2,0)$ & $(2,2,0)$ & $104$\\
$(1,0,3)$ & $(0,0,3)$ & $10$\\
$(2,1,3)$ & $(2,1,2)$ & $97$\\
$(3,1,2)$ & $(2,1,2)$ & $99$\\
$(2,3,3)$ & $(1,3,3)$ & $79$\\
$(1,3,0)$ & $(0,3,0)$ & $34$\\
$(2,0,3)$ & $(1,0,3)$ & $50$\\
$(1,1,3)$ & $(2,1,3)$ & $61$\\
$(3,1,3)$ & $(2,1,3)$ & $101$\\
$(4,1,2)$ & $(3,1,2)$ & $139$\\
$(3,3,3)$ & $(2,3,3)$ & $119$\\
$(2,3,0)$ & $(1,3,0)$ & $74$\\
$(3,0,3)$ & $(2,0,3)$ & $90$\\
$(1,2,3)$ & $(1,1,3)$ & $60$\\
$(4,1,3)$ & $(3,1,3)$ & $141$\\
$(4,2,2)$ & $(4,1,2)$ & $172$\\
$(4,3,3)$ & $(3,3,3)$ & $159$\\
$(3,3,0)$ & $(2,3,0)$ & $114$\\
$(4,0,3)$ & $(3,0,3)$ & $130$\\
$(0,2,3)$ & $(1,2,3)$ & $32$\\
$(2,2,3)$ & $(1,2,3)$ & $72$\\
$(4,2,1)$ & $(4,2,2)$ & $176$\\
$(4,3,0)$ & $(3,3,0)$ & $154$\\
$(3,2,3)$ & $(2,2,3)$ & $112$\\
$(4,3,1)$ & $(4,2,1)$ & $177$\\
$(4,2,3)$ & $(3,2,3)$ & $152$\\
\bottomrule
\end{longtable}
Thus all 80 vertices are reached by 79 selected edges, proving that the selected graph is a tree.  Removing the 15 retained chords partitions it into the following 16 components; the left entry is the unique terminal in the component.
\begin{longtable}{@{}p{.13\textwidth}p{.79\textwidth}@{}}
\toprule terminal & complete vertex set\\ \midrule
\endfirsthead \toprule terminal & complete vertex set\\ \midrule \endhead
$(4,0,0)$ & $(0,0,0)$, $(0,1,0)$, $(1,0,0)$, $(2,0,0)$, $(3,0,0)$, $(4,0,0)$\\
$(4,0,1)$ & $(0,0,1)$, $(0,1,1)$, $(0,2,1)$, $(0,3,1)$, $(1,0,1)$, $(1,1,1)$, $(1,2,1)$, $(1,3,1)$, $(2,0,1)$, $(2,1,1)$, $(2,2,1)$, $(2,3,1)$, $(3,0,1)$, $(3,1,1)$, $(3,2,1)$, $(3,3,1)$, $(4,0,1)$\\
$(4,0,2)$ & $(4,0,2)$\\
$(4,0,3)$ & $(0,0,3)$, $(0,1,3)$, $(1,0,3)$, $(2,0,3)$, $(3,0,3)$, $(4,0,3)$\\
$(4,1,0)$ & $(1,1,0)$, $(2,1,0)$, $(3,1,0)$, $(4,1,0)$\\
$(4,1,1)$ & $(4,1,1)$\\
$(4,1,2)$ & $(0,0,2)$, $(0,1,2)$, $(1,0,2)$, $(1,1,2)$, $(2,0,2)$, $(2,1,2)$, $(3,0,2)$, $(3,1,2)$, $(4,1,2)$\\
$(4,1,3)$ & $(1,1,3)$, $(2,1,3)$, $(3,1,3)$, $(4,1,3)$\\
$(4,2,0)$ & $(4,2,0)$\\
$(4,2,1)$ & $(4,2,1)$\\
$(4,2,2)$ & $(4,2,2)$\\
$(4,2,3)$ & $(0,2,3)$, $(1,2,3)$, $(2,2,3)$, $(3,2,3)$, $(4,2,3)$\\
$(4,3,0)$ & $(0,2,0)$, $(0,3,0)$, $(1,2,0)$, $(1,3,0)$, $(2,2,0)$, $(2,3,0)$, $(3,2,0)$, $(3,3,0)$, $(4,3,0)$\\
$(4,3,1)$ & $(4,3,1)$\\
$(4,3,2)$ & $(1,3,2)$, $(2,3,2)$, $(3,3,2)$, $(4,3,2)$\\
$(4,3,3)$ & $(0,2,2)$, $(0,3,2)$, $(0,3,3)$, $(1,2,2)$, $(1,3,3)$, $(2,2,2)$, $(2,3,3)$, $(3,2,2)$, $(3,3,3)$, $(4,3,3)$\\
\bottomrule
\end{longtable}
Finally delete from the face dual the duals of the 79 gauge-tree edges and the 15 retained chords.  The remainder has 90 face vertices.  The following 89 parent incidences reach every face from $O$; the last column is the primal grid edge crossed by the dual parent edge.
\begin{longtable}{@{}rrr@{}}
\toprule child face & parent face & primal edge index\\ \midrule
\endfirsthead \toprule child face & parent face & primal edge index\\ \midrule \endhead
$(0;0,2,1)$ & $O$ & $35$\\
$(1;0,3,0)$ & $O$ & $33$\\
$(1;0,3,2)$ & $O$ & $37$\\
$(0;0,1,1)$ & $(0;0,2,1)$ & $25$\\
$(1;1,3,0)$ & $(1;0,3,0)$ & $73$\\
$(1;1,3,2)$ & $(1;0,3,2)$ & $77$\\
$(0;0,0,1)$ & $(0;0,1,1)$ & $14$\\
$(1;2,3,0)$ & $(1;1,3,0)$ & $113$\\
$(1;2,3,2)$ & $(1;1,3,2)$ & $117$\\
$(1;3,3,0)$ & $(1;2,3,0)$ & $153$\\
$(1;3,3,2)$ & $(1;2,3,2)$ & $157$\\
$(0;4,2,0)$ & $(1;3,3,0)$ & $181$\\
$(0;4,2,2)$ & $(1;3,3,2)$ & $183$\\
$(0;4,1,0)$ & $(0;4,2,0)$ & $174$\\
$(2;3,2,0)$ & $(0;4,2,0)$ & $175$\\
$(0;4,1,2)$ & $(0;4,2,2)$ & $178$\\
$(0;4,2,1)$ & $(0;4,2,2)$ & $179$\\
$(2;3,2,3)$ & $(0;4,2,2)$ & $180$\\
$(0;4,0,0)$ & $(0;4,1,0)$ & $167$\\
$(2;3,1,1)$ & $(0;4,1,0)$ & $170$\\
$(2;2,2,0)$ & $(2;3,2,0)$ & $143$\\
$(0;4,0,2)$ & $(0;4,1,2)$ & $171$\\
$(2;3,1,3)$ & $(0;4,1,2)$ & $173$\\
$(1;3,3,1)$ & $(0;4,2,1)$ & $182$\\
$(2;2,2,3)$ & $(2;3,2,3)$ & $151$\\
$(1;3,0,0)$ & $(0;4,0,0)$ & $160$\\
$(2;3,0,0)$ & $(0;4,0,0)$ & $161$\\
$(0;3,1,0)$ & $(2;3,1,1)$ & $135$\\
$(2;1,2,0)$ & $(2;2,2,0)$ & $103$\\
$(0;4,0,1)$ & $(0;4,0,2)$ & $165$\\
$(1;3,0,2)$ & $(0;4,0,2)$ & $164$\\
$(2;3,0,3)$ & $(0;4,0,2)$ & $166$\\
$(0;3,1,2)$ & $(2;3,1,3)$ & $140$\\
$(0;3,2,1)$ & $(1;3,3,1)$ & $155$\\
$(2;1,2,3)$ & $(2;2,2,3)$ & $111$\\
$(1;2,0,0)$ & $(1;3,0,0)$ & $120$\\
$(2;2,0,0)$ & $(2;3,0,0)$ & $121$\\
$(1;2,1,0)$ & $(0;3,1,0)$ & $131$\\
$(1;2,2,0)$ & $(0;3,1,0)$ & $142$\\
$(2;3,1,0)$ & $(0;3,1,0)$ & $132$\\
$(2;0,2,0)$ & $(2;1,2,0)$ & $63$\\
$(1;3,1,1)$ & $(0;4,0,1)$ & $169$\\
$(1;2,0,2)$ & $(1;3,0,2)$ & $126$\\
$(2;2,0,3)$ & $(2;3,0,3)$ & $129$\\
$(1;2,1,2)$ & $(0;3,1,2)$ & $137$\\
$(1;2,2,2)$ & $(0;3,1,2)$ & $148$\\
$(2;3,1,2)$ & $(0;3,1,2)$ & $138$\\
$(1;3,2,1)$ & $(0;3,2,1)$ & $145$\\
$(2;2,2,1)$ & $(0;3,2,1)$ & $146$\\
$(2;2,2,2)$ & $(0;3,2,1)$ & $149$\\
$(2;0,2,3)$ & $(2;1,2,3)$ & $71$\\
$(1;1,0,0)$ & $(1;2,0,0)$ & $80$\\
$(2;1,0,0)$ & $(2;2,0,0)$ & $81$\\
$(0;2,1,0)$ & $(1;2,2,0)$ & $102$\\
$(0;3,0,1)$ & $(1;3,1,1)$ & $134$\\
$(1;1,0,2)$ & $(1;2,0,2)$ & $86$\\
$(2;1,0,3)$ & $(2;2,0,3)$ & $89$\\
$(0;2,1,2)$ & $(1;2,2,2)$ & $108$\\
$(0;2,2,1)$ & $(2;2,2,1)$ & $106$\\
$(1;0,0,0)$ & $(1;1,0,0)$ & $40$\\
$(2;0,0,0)$ & $(2;1,0,0)$ & $41$\\
$(2;1,1,0)$ & $(0;2,1,0)$ & $92$\\
$(2;1,1,1)$ & $(0;2,1,0)$ & $95$\\
$(1;3,0,1)$ & $(0;3,0,1)$ & $123$\\
$(2;2,0,1)$ & $(0;3,0,1)$ & $124$\\
$(2;2,0,2)$ & $(0;3,0,1)$ & $127$\\
$(1;0,0,2)$ & $(1;1,0,2)$ & $46$\\
$(2;0,0,3)$ & $(2;1,0,3)$ & $49$\\
$(2;1,1,2)$ & $(0;2,1,2)$ & $98$\\
$(2;1,1,3)$ & $(0;2,1,2)$ & $100$\\
$(1;1,2,1)$ & $(0;2,2,1)$ & $105$\\
$(1;1,3,1)$ & $(0;2,2,1)$ & $115$\\
$(0;1,1,0)$ & $(2;1,1,1)$ & $55$\\
$(0;2,0,1)$ & $(2;2,0,1)$ & $84$\\
$(0;1,1,2)$ & $(2;1,1,2)$ & $58$\\
$(0;1,2,1)$ & $(1;1,3,1)$ & $75$\\
$(1;0,1,0)$ & $(0;1,1,0)$ & $51$\\
$(1;0,2,0)$ & $(0;1,1,0)$ & $62$\\
$(1;1,0,1)$ & $(0;2,0,1)$ & $83$\\
$(1;1,1,1)$ & $(0;2,0,1)$ & $94$\\
$(1;0,1,2)$ & $(0;1,1,2)$ & $57$\\
$(1;0,2,2)$ & $(0;1,1,2)$ & $68$\\
$(2;0,2,1)$ & $(0;1,2,1)$ & $66$\\
$(2;0,2,2)$ & $(0;1,2,1)$ & $69$\\
$(0;0,1,0)$ & $(1;0,2,0)$ & $22$\\
$(0;1,0,1)$ & $(1;1,1,1)$ & $54$\\
$(0;0,1,2)$ & $(1;0,2,2)$ & $28$\\
$(2;0,0,1)$ & $(0;1,0,1)$ & $44$\\
$(2;0,0,2)$ & $(0;1,0,1)$ & $47$\\
\bottomrule
\end{longtable}
The table is a spanning tree of the retained dual remainder and therefore proves its connectivity.  Together, these three printed witnesses establish the tree, one-terminal-per-component, and dual connectivity assertions for the opposite width-four base.
For the exceptional width-four relation, remove also edge 37 from the retained dual remainder.  The component containing $O$ is the following complete face set:
\begin{quote}\small $O$, $(1;0,0,0)$, $(2;0,0,0)$, $(0;0,0,1)$, $(1;0,1,0)$, $(0;0,1,0)$, $(0;0,1,1)$, $(1;0,2,0)$, $(2;0,2,0)$, $(0;0,2,1)$, $(1;0,3,0)$, $(1;1,0,0)$, $(2;1,0,0)$, $(2;1,1,0)$, $(0;1,1,0)$, $(2;1,1,1)$, $(2;1,2,0)$, $(1;1,3,0)$, $(1;2,0,0)$, $(2;2,0,0)$, $(1;2,1,0)$, $(0;2,1,0)$, $(1;2,2,0)$, $(2;2,2,0)$, $(1;2,3,0)$, $(1;3,0,0)$, $(2;3,0,0)$, $(2;3,1,0)$, $(0;3,1,0)$, $(2;3,1,1)$, $(2;3,2,0)$, $(1;3,3,0)$, $(0;4,0,0)$, $(0;4,1,0)$, $(0;4,2,0)$.\end{quote}
Its non-gauge chord boundary consists exactly of primal edge indices $37,163,177$, namely $X=e_z(0,3,2)$, $e_y(4,0,1)$, and $e_y(4,2,1)$.  Its face-boundary relation is therefore $h_X=h_{e_y(4,0,1)}+h_{e_y(4,2,1)}$.
For the terminal relation, deleting edge 37 leaves terminal set $\{(4,3,3)\}$ away from the root; deleting edge 65 ($e_z(1,2,1)$) leaves $\{(4,3,2),(4,3,3)\}$; and deleting edge 109 ($e_y(2,2,2)$) leaves $\{(4,3,2)\}$.  Symmetric difference therefore gives $u_X=u_{e_z(1,2,1)}+u_{e_y(2,2,2)}$.
